\section*{Bab \ref{ch:g11:dissolution} --- Melarutkan Padatan Ion dan Padatan Molekul}

\begin{solution}{exo:g11:dissolution:1}
\ce{NaCl(s) -> Na+(aq) + Cl-(aq)};
\ce{CaCl2(s) -> Ca^{2+}(aq) + 2Cl-(aq)};
\ce{Na2SO4(s) -> 2Na+(aq) + SO4^{2-}(aq)};
\ce{AlCl3(s) -> Al^{3+}(aq) + 3Cl-(aq)}.
\end{solution}

\begin{solution}{exo:g11:dissolution:2}
$[\ce{Na+}] = 2 \times 0.050 = \qty{0.10}{mol/L}$;
$[\ce{SO4^{2-}}] = \qty{0.050}{mol/L}$.
\end{solution}

\begin{solution}{exo:g11:dissolution:3}
\omterm{def:g11:dissolution:solvation}{Hidrasi} adalah dikelilinginya \omterm{def:g9:ions:ion}{ion} itu oleh \omterm{def:g7:atoms-and-molecules:molecule}{molekul-molekul} air yang
terikat padanya. Ujung oksigen ($\delta^-$) mengarah ke \omterm{def:g9:ions:ion}{kation}, ujung
hidrogen ($\delta^+$) ke \omterm{def:g9:ions:ion}{anion}: muatan yang berlawanan saling tarik.
\end{solution}

\begin{solution}{exo:g11:dissolution:4}
$M(\ce{AlCl3}) = 27.0 + 3 \times 35.5 = \qty{133.5}{g/mol}$;
$n = 2.67 / 133.5 = \qty{0.0200}{mol}$; $c = 0.0200 / 0.2000 =
\qty{0.100}{mol/L}$; $[\ce{Al^{3+}}] = \qty{0.100}{mol/L}$,
$[\ce{Cl-}] = \qty{0.300}{mol/L}$.
\end{solution}

\begin{solution}{exo:g11:dissolution:5}
Kepala \ce{-COO-} bersifat ionik, ditarik oleh air: \omterm{def:g11:dissolution:hydrophilic}{hidrofilik}. Ekor
dari 17 \omterm{def:g7:atoms-and-molecules:atom}{atom} karbon, dengan hanya ikatan \ce{C-C} dan ikatan \ce{C-H}
yang hampir nonpolar, bersifat nonpolar: \omterm{def:g11:dissolution:hydrophilic}{hidrofobik}.
\end{solution}

\begin{solution}{exo:g11:dissolution:6}
Garam: air (padatan \omterm{def:g9:ions:ion}{ion}, \omterm{def:g6:solutions-and-solubility:solute}{pelarut} polar). Lilin: sikloheksana (rantai
nonpolar). Gula: air (gugus \ce{O-H}-nya membentuk \omterm{def:g11:polarity-and-cohesion:hydrogen-bond}{ikatan hidrogen}
dengan air). Iodin: sikloheksana (\omterm{def:g11:polarity-and-cohesion:polar-molecule}{molekul nonpolar}).
\end{solution}

\begin{solution}{exo:g11:dissolution:7}
Gugus \ce{O-H} pada etanol membentuk \omterm{def:g11:polarity-and-cohesion:hydrogen-bond}{ikatan hidrogen} dengan \omterm{def:g7:atoms-and-molecules:molecule}{molekul-molekul}
air, yang menggantikan ikatan yang diputusnya: etanol masuk ke dalam
air. Heksana, yang nonpolar, tidak dapat membentuk ikatan seperti itu;
\omterm{def:g7:atoms-and-molecules:molecule}{molekul-molekul} air, yang terikat satu sama lain oleh \omterm{def:g11:polarity-and-cohesion:hydrogen-bond}{ikatan hidrogen},
tidak membiarkannya masuk.
\end{solution}

\begin{solution}{exo:g11:dissolution:8}
Ekor-ekor \omterm{def:g11:dissolution:hydrophilic}{hidrofobik} menghindari air dan \omterm{def:g2:mixing-and-dissolving:dissolve}{larut} di dalam lemak, di
tengah; kepala-kepala bermuatan ditarik oleh air dan tetap di luar.
Semua \omterm{def:g11:dissolution:amphiphilic}{misel} membawa muatan negatif di permukaannya, dan muatan sejenis
saling tolak: \omterm{def:g11:dissolution:amphiphilic}{misel-misel} itu tetap terpisah.
\end{solution}

\begin{solution}{exo:g11:dissolution:9}
Sikloheksana: \omterm{def:g6:solutions-and-solubility:solute}{pelarut} ini \omterm{def:g6:solutions-and-solubility:miscible}{tidak dapat bercampur} dengan air, sehingga
membentuk lapisan tersendiri yang dapat dikeluarkan, dan zat harum itu
(yang nonpolar) mudah \omterm{def:g2:mixing-and-dissolving:dissolve}{larut} di dalamnya. Etanol \omterm{def:g6:solutions-and-solubility:miscible}{dapat bercampur} dengan
air: tidak terbentuk lapisan kedua, dan tidak ada yang dapat dipisahkan.
\end{solution}

\begin{solution}{exo:g11:dissolution:10}
Massa jenis sikloheksana lebih kecil daripada air. Lapisan bawah, yang
berair, dikeluarkan melalui keran. Iodin berada di lapisan sikloheksana
di atas, yang tetap di dalam corong.
\end{solution}

\begin{solution}{exo:g11:dissolution:11}
Tidak ada yang berubah: tembaga sulfat bersifat ionik dan tidak \omterm{def:g2:mixing-and-dissolving:dissolve}{larut}
dalam sikloheksana yang nonpolar. Airnya tetap biru dan sikloheksananya
tak berwarna.
\end{solution}

\begin{solution}{exo:g11:dissolution:12}
$[\ce{Cl-}] = 2c$, jadi $c = \qty{0.150}{mol/L}$;
$n = 0.150 \times 0.2500 = \qty{0.0375}{mol}$;
$m = 0.0375 \times 111.1 = \qty{4.17}{g}$.
\end{solution}

\begin{solution}{exo:g11:dissolution:13}
\omterm{def:g9:ions:ion}{Ion-ion} kalsium dan magnesium mengendapkan \omterm{def:g9:ions:ion}{ion-ion} stearat sebagai kerak
sabun: sabun habis terpakai sebelum dapat membentuk \omterm{def:g11:dissolution:amphiphilic}{misel}, dan sulit
berbusa. Kalsium saja: \qty{0.010}{mol} \ce{Ca^{2+}} mengambil
\qty{0.020}{mol} stearat, yaitu $0.020 \times 306.0 =
\qty{6.1}{g}$ natrium stearat per liter.
\end{solution}

\begin{solution}{exo:g11:dissolution:14}
Volume \qty{0.2000}{L}. \ce{Na+}: $0.020 / 0.2000 = \qty{0.10}{mol/L}$;
\ce{Ca^{2+}}: $0.010 / 0.2000 = \qty{0.050}{mol/L}$; \ce{Cl-}:
$(0.020 + 0.020) / 0.2000 = \qty{0.20}{mol/L}$. Muatan positif:
$0.10 + 2 \times 0.050 = 0.20$; negatif: $0.20$. Netral.
\end{solution}

\begin{solution}{exo:g11:dissolution:15}
Satu liter heptana bermassa \qty{0.68}{kg} dan dapat menampung
$17 \times 0.68 = \qty{12}{g}$ iodin, sekitar 40 kali lebih banyak
daripada satu liter air. Jika dikocok bersama, hampir semua iodin
berpindah ke dalam heptana: \omterm{def:g10:chemical-species:extraction}{ekstraksinya} efisien.
\end{solution}

\begin{solution}{pb:g11:dissolution:1}
\textbf{1.} Dari batuan yang dilalui air itu (batu kapur, dolomit,
gips), yang sedikit \omterm{def:g2:mixing-and-dissolving:dissolve}{larut}.

\textbf{2.} $[\ce{Ca^{2+}}] = 0.120 / 40.1 = \qty{2.99e-3}{mol/L}$;
$[\ce{Mg^{2+}}] = 0.024 / 24.3 = \qty{9.9e-4}{mol/L}$.

\textbf{3.} \ce{Ca(HCO3)2 -> Ca^{2+}(aq) + 2HCO3-(aq)}; jadi
$[\ce{HCO3-}] = 2 \times \num{2.99e-3} = \qty{5.99e-3}{mol/L}$.

\textbf{4.} $\num{2.99e-3} \times 150 = \qty{0.449}{mol}$.

\textbf{5.} $18 \times 12.0 + 35 \times 1.0 + 2 \times 16.0 + 23.0 =
\qty{306.0}{g/mol}$.

\textbf{6.} \ce{C17H35COONa(s) -> C17H35COO-(aq) + Na+(aq)}.

\textbf{7.} Kepala \ce{-COO-}: \omterm{def:g11:dissolution:hydrophilic}{hidrofilik}; ekor \ce{C17H35-}:
\omterm{def:g11:dissolution:hydrophilic}{hidrofobik}. \omterm{def:g9:ions:ion}{Ion} itu memiliki kedua jenis bagian: \omterm{def:g11:dissolution:amphiphilic}{amfifilik}.

\textbf{8.} Sebuah bulatan kecil \omterm{def:g9:ions:ion}{ion-ion} sabun, ekor ke dalam, kepala ke
luar. Ekor-ekornya \omterm{def:g2:mixing-and-dissolving:dissolve}{larut} di dalam lemak, yang pecah menjadi tetes-tetes
yang terbungkus \omterm{def:g11:dissolution:amphiphilic}{misel}; permukaan yang bermuatan menjaga tetes-tetes itu
tetap terpisah di dalam air, yang lalu membilasnya.

\textbf{9.} Sabun harus membentuk \omterm{def:g9:ions:ion}{ion} dan \omterm{def:g11:dissolution:amphiphilic}{misel}, yang memerlukan air di
sekeliling kepala-kepalanya; dan lemak dibawa pergi oleh air bilasan.
Dalam minyak, lemak itu akan tetap terlarut di atas kulit.

\textbf{10.} \ce{2C17H35COO- + Ca^{2+} -> Ca(C17H35COO)2}: muatan
$2 \times (-1) + 2 = 0$ di kiri, $0$ di kanan.

\textbf{11.} \ce{C17H35COO-}: berlebih $- 2x$; \ce{Ca^{2+}}:
$\num{2.99e-3} - x$; kalsium stearat: $x$. \omterm{def:g9:ions:ion}{Ion-ion} kalsium adalah
pembatas: $x_{\max} = \qty{2.99e-3}{mol}$.

\textbf{12.} $2 \times \num{2.99e-3} = \qty{5.99e-3}{mol}$ per liter.

\textbf{13.} $M = 36 \times 12.0 + 70 \times 1.0 + 4 \times 16.0 + 40.1 =
\qty{606.1}{g/mol}$; $\num{2.99e-3} \times 606.1 = \qty{1.81}{g}$ kerak
sabun per liter.

\textbf{14.} $2 \times \num{9.9e-4} = \qty{1.98e-3}{mol}$ per liter.

\textbf{15.} $2 \times 0.449 = \qty{0.898}{mol}$.

\textbf{16.} $0.898 \times 306.0 = \qty{275}{g}$.

\textbf{17.} Magnesium: $\num{9.9e-4} \times 150 = \qty{0.148}{mol}$,
yang mengambil \qty{0.296}{mol} stearat, \qty{91}{g} sabun; seluruhnya
sekitar $275 + 91 = \qty{366}{g}$.

\textbf{18.} $275 / 100 \approx 2.7$ batang.

\textbf{19.} Kalsium dalam satu bak mandi membuang sekitar
\qty{0.27}{kg} sabun sebagai kerak sabun.
\end{solution}
